% Laget automatisk av elevark/lag_elevark.py fra sim/bevegelsesgrafer.html. Endringer her blir overskrevet.
\documentclass[11pt]{article}
\usepackage{elevark}
\begin{document}
\elevarkhode{Fysikk 1 · Rettlinjet bevegelse}{Bevegelsesgrafer}{https://lektorbo.no/sim/bevegelsesgrafer.html}{Simuleringen}

En bevegelse kan beskrives med tre grafer: posisjon, fart og akselerasjon mot tid. Bygg en bevegelse av opptil fire deler med konstant akselerasjon, og se hvordan grafene henger sammen.

\section*{Oppgaver}

\begin{oppgaveliste}
\oppgave \fase{gjett}{Gjett}Velg «Konstant fart». Hvordan ser de tre grafene ut? Hvordan ser du farten på posisjonsgrafen? Hvordan ligger merkene for hvert halve sekund?
\svar
\oppgave \fase{test}{Test}Velg «Fra ro med konstant akselerasjon». Hvordan endrer avstanden mellom merkene seg? Hvorfor er posisjonsgrafen krum?
\svar
\oppgave \fase{test}{Regn}Velg «Bil som akselererer og bremser». Bruk tallene i tabellen til å regne ut toppfarten og hvor langt bilen kjører i hver del. Slå på «Areal under fartsgrafen» og sjekk.
\svar
\oppgave \fase{forklar}{Forklar}Velg «Bremser, snur og kjører tilbake». Når står bilen stille? Hva er farten og akselerasjonen akkurat da? Forklar forskjellen på tilbakelagt strekning og forflytning når bevegelsen er ferdig.
\svar
\oppgave \fase{test}{Lag}Lag en bevegelse der bilen starter i ro, kjører 30~m fram og stopper. Prøv å bruke akkurat tre deler med like lang varighet.
\svar
\oppgave \fase{test}{Spill}Spill «Følg grafen». Hva må du gjøre med farten når den stiplede grafen er bratt? Når den heller nedover?
\svar
\oppgave \fase{utfordring}{Utfordring}Kan en bil ha fart null og samtidig akselerasjon som ikke er null? Kan den ha negativ fart og positiv akselerasjon? Lag eksempler i simuleringen.
\svar
\end{oppgaveliste}

\end{document}
